\documentclass[10pt,a4paper]{amsart}
\usepackage[margin=19mm]{geometry}
\usepackage{amsmath,amssymb,mathtools}
\usepackage[scr=rsfs,cal=euler]{mathalfa}
\usepackage{xfrac}
\usepackage[unicode,hidelinks,bookmarks=false]{hyperref}
\newcommand{\ie}{i.e.\ }
\newcommand{\cf}{cf.\ }
\newcommand{\resp}{respectively\ }
\newcommand{\Subsection}{\subsection}
\providecommand{\nobreakdash}{\nobreak\hbox{-}}
\setlength{\emergencystretch}{2em}
\multlinegap0pt
\usepackage{bm}
\usepackage{bbm}
\usepackage{dsfont}
\usepackage{xspace}
\usepackage{cancel}
\usepackage{mathtools}
\usepackage{amsthm}
\makeatletter
\@ifundefined{thm}{\newtheorem{thm}{Theorem}}{}
\makeatother
\def\sC{{\mathfrak C}}
\def\uphar{{\upharpoonright\,}}
\newcommand{\cD}{\mathcal{D}\,}
\newcommand{\cH}{\mathcal{H}\,}
\newcommand{\cJ}{\mathcal{J}\,}
\newcommand{\cM}{\mathcal{M}\,}
\newcommand{\cN}{\mathcal{N}}
\newcommand{\cQ}{\mathcal{Q}}
\newcommand{\cR}{\mathcal{R}}
\newcommand{\cS}{\mathcal{S}}
\newcommand{\gA}{\mathfrak{A}\,}
\newcommand{\gC}{\mathfrak{C}}
\newcommand{\gT}{\mathfrak{T}}
\newcommand{\ii}{\mathrm{i}}
\title[On $C$-Symmetric and $C$-Self-adjoint Unbounded Operators on]{On $C$-Symmetric and $C$-Self-adjoint Unbounded Operators on Hilbert Space}
\author{Yury Arlinskii, Konrad Schm\"udgen}
\date{Source: arXiv:2507.01847v2 (CC BY 4.0)}
\begin{document}
\maketitle

\noindent\emph{Source and licence.} On $C$-Symmetric and $C$-Self-adjoint Unbounded Operators on Hilbert Space. Version arXiv:2507.01847v2 (CC BY 4.0), distributed under the Creative Commons Attribution 4.0 International licence, as recorded in the arXiv licence metadata for this exact version and shown on the abstract page https://arxiv.org/abs/2507.01847v2. Full rights notice: licenses/arxiv-2507.01847-CC-BY-4.0.txt.\par\smallskip

\noindent\textit{Editorial scope.} Counted unit: the Thm whose source label is \texttt{rev02ba} in the article (source0.tex lines 794--818, source label \texttt{rev02ba}), delivered together with the article's own complete original proof of it (source0.tex lines 819--892, one \texttt{proof} environment of 74 lines). No mathematics is written, repaired or re-derived: statement and proof are the article's own text line for line. Required context retained uncounted: 22agua (lines 260--262); matrixtbt1 (lines 434--440); gua30a (lines 578--599); 21abgua (lines 584--589); 21bcgua (lines 591--597); defics (lines 632--634); matrixtbt3 (lines 679--687); conditionCU (lines 693--694); rev02a (lines 770--777); ver02ba (lines 779--787). Every definition, display and label of those lines is retained unchanged. Omitted from this package and not redistributed: 1--259, 263--433, 441--577, 590--590, 598--631, 635--678, 688--692, 695--769, 778--778, 788--793, 893--1396. Nothing inside a retained excerpt is omitted or altered: every retained body is delivered exactly as the article's lines are written, so no omission receipt of any kind accompanies this package. Citations: the retained excerpts carry no citation command at all, so no bibliography entry of the article is transcribed and no reference is redistributed. Reference adaptation: none was needed and none was made; every reference inside the delivered unit resolves inside it, so no unresolved reference or citation is produced. Printed slips kept literal: no printed word, symbol, hypothesis or number of the counted statement or of its proof was altered. Layout and class: the article's own class is \texttt{amsart} with packages including amssymb, amsthm, amsfonts, amsmath, color; the wrapper is the lane's pinned \texttt{amsart} editorial wrapper at 10pt on A4, which restates the theorem-like environments the retained text needs and adds the article's own macro definitions that the retained text needs (\texttt{\textbackslash cD}, \texttt{\textbackslash cH}, \texttt{\textbackslash cJ}, \texttt{\textbackslash cM}, \texttt{\textbackslash cN}, \texttt{\textbackslash cQ}, \texttt{\textbackslash cR}, \texttt{\textbackslash cS}, \texttt{\textbackslash gA}, \texttt{\textbackslash gC}, \texttt{\textbackslash gT}, \texttt{\textbackslash ii}, \texttt{\textbackslash sC}, \texttt{\textbackslash uphar}), with extra wrapper packages (bm, bbm, dsfont, xspace, cancel, mathtools). The delivered numbering is the unit's own. Assets: no retained span contains a figure environment, a graphics-inclusion command, a PDF-page inclusion, a table or a TikZ picture; no image file is copied into this package, the package declares no asset dependency and no image file is redistributed with it. Licence: version arXiv:2507.01847v2 of this article is distributed under the Creative Commons Attribution 4.0 International licence, as recorded in the arXiv licence metadata for this exact version and shown on the abstract page https://arxiv.org/abs/2507.01847v2. A full rights notice is delivered at licenses/arxiv-2507.01847-CC-BY-4.0.txt. Characters: every retained excerpt is pure ASCII, so the delivered engine, XeLaTeX with its default Latin Modern text font, reports no missing character.
\begin{equation}\label{22agua}
\langle f,g\rangle_{\cH_T}=\langle f,g\rangle +\langle Tf,Tg\rangle, \quad f,g\in \cD(T).
\end{equation}

\begin{gather}\label{matrixtbt1}
\gA=\left(
\begin{array}{ll}
~~~ 0 & A \\
CA\,C & 0
\end{array}\right).
\end{gather}

\begin{thm}\label{gua30a}
Let $A$ be a densely defined closed $C$-symmetric operator in $\cH$. There is a bijective correspondence between all $C$-self-adjoint extensions of $A$ and all conjugations $\cJ$ of the subspace $\cN((A^*C)^2+I)$ w.r.t. the inner product in $\cH_{CA^*C}$ (see \eqref{22agua}), satisfying the anticommutation relation
\begin{equation}\label{30aagua}
A^*C\cJ h=-\cJ A^*Ch, \quad h\in \cN((A^*C)^2+I).
\end{equation}
 The correspondence is given by the formulas
 \begin{equation}\label{21abgua}
\left\{\begin{array}{l}  \cD(A_{\cJ})=\cD(A)\dot+(I-A^*C \cJ)\cN((A^*C)^2+I),\\[2mm]
A_{\cJ} (f+(I-A^*C \cJ)h)=Af+C(A^*C+\cJ)h,\\[2mm]
 f\in\cD(A),\; h\in\cN( (A^*C)^2+I).
\end{array}\right.
  \end{equation}
  The adjoint operator $(A_{\cJ})^*=CA_\cJ C$ is of the form
\begin{equation}\label{21bcgua}
\left\{ \begin{array}{l}
\cD((A_\cJ)^*)=C\cD(A)\dot+C(I-A^*C\cJ)\cN((A^*C)^2+I),\\[2mm] %%\\
(A_\cJ)^*(Cf+C(I-A^*C\cJ)g)=CAf+(A^*C+\cJ)g,\\[2mm]
f\in\cD(A),\; g\in\cN( (A^*C)^2+I).
\end{array}\right.
\end{equation}

\end{thm}

\begin{align}\label{defics}
\cS:=A^*C\uphar\cN((A^*C)^2+I).
\end{align}

\begin{equation}
\label{matrixtbt3}
\left\{\begin{array}{l}\gT=
\begin{pmatrix}
 0 & \widetilde A\cr C\widetilde A C & 0
\end{pmatrix},\\[2mm]
\cD(\gT)=\left\{(Cf_{\widetilde A},g_{\widetilde A}): f_{\widetilde A},\; g_{\widetilde A}\in\cD(\widetilde A)\right\}=C\cD(\widetilde A)\oplus\cD(\widetilde A).
\end{array}\right.
\end{equation}

\begin{align}\label{conditionCU}\gC\,U \gC\, U=I.
\end{align}

\begin{equation}\label{rev02a}
\begin{array}{l}
\cD(\gT)=C\cD(\widetilde A)\oplus\cD(\widetilde A)%%\\[2mm]
=C\cD(A)\oplus\cD(A)\dot+(I-U)\cN_+\\[2mm]
\left\{(Cf_A-\ii C(\cS-\cJ)h,g_A+(I-\cS \cJ) h),\; %%\\
f_A,g_A\in\cD(A),\;h\in\cN(\cS^2+I)\right\}
\end{array}
\end{equation}

\begin{equation}\label{ver02ba}
\begin{array}{l}
\gT\left((Cf_A-\ii C(\cS-\cJ)h,g_A+(I-\cS \cJ) h)\right)\\
=\left(Ag_A,CAf_A\right)
+\ii\left(-\ii C(\cS+\cJ)h,(I+\cS \cJ) h\right)\\[2mm]
=\left(Ag_A+C(\cS+\cJ)h, CAf_A+\ii(I+\cS \cJ) h\right),\;%%\\[2mm]
 f_A,g_A\in\cD(A),\;h\in\cN(\cS^2+I).
\end{array}
\end{equation}

\begin{thm}\label{rev02ba}
Suppose that $A$ is a densely defined closed $C$-symmetric operator on $\cH$ and  $\gA$ denotes the operator  given by the block matrix \eqref{matrixtbt1}. Let $\widetilde \gA$ be a self-adjoint and $\sC$-self-adjoint extension of  $ \gA$ and let $U$ be the unitary operator of $\cN_+$ on $\cN_-$, which determines the extension $\widetilde \gA$ via von Neumann formulas. Then
the following are equivalent:
\begin{enumerate}
%%\def\labelenumi{\rm (\roman{enumi})}
\item[ (a)]
The operator $\widetilde \gA$ can be represented by the operator matrix
\begin{equation} \label{ver03ea}
%%%\begin{array}{l}
\left\{ \begin{array}{l}\widetilde\gA=\begin{pmatrix}0&\widetilde F\cr\widetilde G &0\end{pmatrix},\\[2mm]
\cD(\widetilde\gA)=%%\left\{(g,f):f\in\cD(\widetilde F), g\in\widetilde \cD(G)\right\}=
\cD(\widetilde G)\oplus \cD(\widetilde F).
\end{array}\right.
\end{equation}
\item [ (b)]
The operator $U$ takes the form
\begin{equation}\label{rev04a}
\left\{\begin{array}{l}U(-\ii CA^*C h, h)=(-\ii  C\cJ h, A^*C\cJ h),\; h\in\cN(\cS^2+I),\\[2mm]
\cJ \mbox{is a conjugation in} \;\; \cN((A^*C)^2+I)\; \mbox{and}\;\;A^*C\cJ=-\cJ A^*C.
\end{array}\right.
\end{equation}
\end{enumerate}
Moreover, if condition \eqref{rev04a} is fulfilled, then $\widetilde F=A_\cJ$, $\widetilde G=CA_\cJ C$, where $A_\cJ$ is $C$-self-adjoint extension of $A$ of the form
\eqref{21abgua}.
\end{thm}

\begin{proof}
As it has been mentioned in the proof of Theorem \ref{gua30a}, there is a bijective correspondence between all self-adjoint and $\sC$-self-adjoint extensions $\widetilde \gA$ of $\gA$ and  unitary operators $U$ of $\cN_+$ on $\cN_-$ satisfying equation \eqref{conditionCU}. This correspondence is given by the
formulas
\begin{equation}\label{rev04b}
\left\{\begin{array}{l}\cD(\widetilde \gA):=\cD(\gA)\dot{+}(I-U)\cN_+,
\\[2mm]
\widetilde \gA:=\gA^*\lceil \cD(\gA)\Longleftrightarrow \widetilde \gA(f+(I-U)v) =\gA f+\ii \, (I+U)v,\\[2mm]
\textrm{for}~~~ f=(Cf_A,g_A),\,  f_A,g_A\in \cD(A), \,v=(x_+,y_+)\in \cN_+.
\end{array}\right.
\end{equation}
It is already proved in the preceding Theorem \ref{gua30a} that $U$ satisfies \eqref{conditionCU} if and only if $U$ is of the
form
\begin{equation}\label{rev05a}
U(-\ii C\cS h, h)=(-\ii C\cJ h, \cS \cJ h),\; h\in\cN(\cS^2+I),
\end{equation}
where $\cJ$ is a conjugation in $\cN(\cS^2+I)$ and $\cS:=A^*C\uphar \cN((A^*C)^2+I)$, see (\ref{defics}).
Hence the formulas \eqref{rev02a} and \eqref{ver02ba} for the domain and the action of the operator $\widetilde \gA$ are valid.

We also proved in Theorem \ref{gua30a} that
 if $\widetilde A$ is a $C$-self-adjoint extension  of $A$, then the operator matrix  $\widetilde\gA=\begin{pmatrix}
 0 & \widetilde A\cr C\widetilde A C & 0
\end{pmatrix}$ is a self-adjoint and $\sC$-self-adjoint extension of $\gA$ (see \eqref{matrixtbt3}),
$\cS\cJ=-\cJ\cS$ and $\widetilde A$ takes the form \eqref{21abgua},
i.e., we have (b) $\Longrightarrow $ (a) with $\widetilde F=A_\cJ,$ $\widetilde G=CA_\cJ C$.

Now we  prove the implication (a)$\Longrightarrow$ (b).

Let $\widetilde \gA$ be of the form
\eqref{ver03ea}, then
from \eqref{rev04b} and \eqref{rev05a} (see also \eqref{rev02a}) it follows that
\begin{multline*}
\cD(\widetilde\gA)=\cD(\widetilde G)\oplus \cD(\widetilde F)=\\[2mm]
=\left\{(Cf_A-\ii C(\cS-\cJ)h,g_A+(I-\cS \cJ) h):%%\\
\;\;f_A,g_A\in\cD(A),\;h\in\cN(\cS^2+I)\right\}.
\end{multline*}
Note that $$Cf_A-\ii C(\cS-\cJ)h=0\Longleftrightarrow f_A=0,\;(\cS-\cJ)h=0.$$
Hence
$$(\{0\},\cD(\widetilde F))=\left\{(0,g_A+ (I-\cS\cJ)h): g_A\in\cD(A), h\in\cN(\cS-\cJ)\right\}.$$%%\Longleftrightarrow=\left\{f_A\in %%\cD(A)(I-\cS\cJ)\cN(\cS-\cJ)=\cR(I-\cS\cJ).$$
Since $\cS-\cJ=\cS(I+\cS\cJ)$, we get  $\cN(\cS-\cJ)=\cN(I+\cS\cJ)$. Consequently,
\begin{equation}\label{ver03ee}
(I-\cS\cJ)\cN(I+\cS\cJ)=\cR(I-\cS\cJ).
\end{equation}
Because $\cS\cJ$ is a unitary operator, it has a matrix representation of the form
\[%%\begin{equation}\label{rev02e}
\cS\cJ=\begin{pmatrix}  -I_{\cM_-}&0&0\cr
0&I_{\cM_+}&0\cr 0&0&\cQ\end{pmatrix}
\]%%\end{equation}
w.r.t. the orthogonal decomposition
$$ \cN(\cS^2+I)=\cM_-\oplus\cM_+\oplus\cM_0,$$
where $\cM_\mp:=\cN(I\pm\cS\cJ)$,
$\cM_0:=\cN(\cS^2+I)\ominus\left(\cM_-\oplus\cM_+\right)$,  the operator $\cQ$ is unitary in $\cM_0$, and $\cN(I\pm\cQ)=\{0\}.$

It follows that
\[
I-\cS\cJ=\begin{pmatrix}  2I_{\cM_-}&0&0\cr
0&0&0\cr 0&0&I_{\cM_0}-\cQ\end{pmatrix}. %%,\; I+\cS\cJ=\begin{pmatrix}  0&0&0\cr0&2I_{\cM_+}&0\cr 0&0&I_{\cM_0}+\cQ\end{pmatrix}
\]
Hence
\[
\cR(I-\cS\cJ)=\cM_-\oplus\cR(I_{\cM_0}-\cQ).%%=\cM_-.
\]
On the other side, the equality \eqref{ver03ee} gives that $\cR(I-\cS\cJ)=\cM_-$.
So, $\cM_0=\{0\}$ and $\cS\cJ$ is of the form
\[
\cS\cJ=\begin{pmatrix}  -I_{\cM_-}&0\cr
0&I_{\cM_+}\end{pmatrix}.
\]
Consequently, the unitary operator $\cS\cJ$ is self-adjoint in $\cN(\cS^2+I)$, i.e., $\cS\cJ=(\cS\cJ)^{-1}=-\cJ\cS.$
This proves the implication (a)$\Longrightarrow$(b).

Moreover, by Theorem \ref{gua30a} the equalities $\widetilde F=A_\cJ$ and $\widetilde G=C\widetilde A_\cJ C=(A_\cJ)^*$ hold (see \eqref{21abgua}
and \eqref{21bcgua}).

\end{proof}

\end{document}
